%=========== ΛΥΣΗ - 24ο ΘΕΜΑ Α ===========
\providecommand{\theoriaref}[3]{%
\begin{center}
\fcolorbox{maincolor!70!black}{maincolor!8}{%
\begin{minipage}{.88\linewidth}
\textcolor{maincolor!80!black}{\kerkissans{\textbf{#1~\ref{#2}}}}\\[1mm]
#3
\end{minipage}}
\end{center}}
\providecommand{\orismosref}[1]{%
\theoriaref{Ορισμός}{#1}{Η απάντηση δίνεται από τον αντίστοιχο ορισμό.}}
\providecommand{\apodeixiref}[1]{%
\theoriaref{Θεώρημα}{#1}{Η ζητούμενη απόδειξη δίνεται στην αντίστοιχη απόδειξη.}}

\begin{thema}{Α}
\begin{erwthma}
\item \theoriaref{Θεώρημα}{thewrhma:bolzano}{Η διατύπωση και η γεωμετρική ερμηνεία δίνονται στην αντίστοιχη ενότητα.}

\item \orismosref{orismos:praxeis_synarthsewn}

\item \orismosref{orismos:koilh_synarthsh}

\item Οι χαρακτηρισμοί των προτάσεων είναι:
\begin{alist}
\item \textbf{Σωστό}. Αν το $x_0$ είναι εσωτερικό σημείο και θέση τοπικού ακρότατου, τότε είτε $f'(x_0)=0$ είτε η $f'$ δεν ορίζεται στο $x_0$. Άρα το $x_0$ είναι κρίσιμο σημείο.
\item \textbf{Σωστό}. Είναι ο κανόνας παραγώγισης σύνθετης συνάρτησης:
\[ (f(g(x)))'=f'(g(x))\cdot g'(x). \]
\item \textbf{Λάθος}. Η σχέση $f(f^{-1}(x))=x$ ισχύει για κάθε $x$ στο πεδίο ορισμού της $f^{-1}$, δηλαδή για $x\in f(D_f)$, όχι αναγκαστικά για κάθε $x\in D_f$.
\item \textbf{Σωστό}. Για κάθε συνεχή συνάρτηση $f$ ισχύει
\[ \dint_a^a f(x)\d x=0. \]
\item \textbf{Σωστό}. Αν οι συνεχείς $f,g$ δεν είχαν κοινό σημείο, τότε η $f-g$ θα διατηρούσε σταθερό πρόσημο στο $[a,\beta]$, άρα το ολοκλήρωμά της δεν θα ήταν μηδέν. Επομένως υπάρχει $x_0\in[a,\beta]$ με $f(x_0)=g(x_0)$.
\end{alist}

\item Σωστή απάντηση είναι η \textbf{γ}. Από την οριζόντια ασύμπτωτη $y=2$ στο $+\infty$ προκύπτει
\[ \lim_{x\to+\infty}f(x)=2. \]
\end{erwthma}
\end{thema}
%=========== ΤΕΛΟΣ ΛΥΣΗΣ - 24ο ΘΕΜΑ Α ===========
